\chapter*{Danh mục từ viết tắt và thuật ngữ}
\label{Appendix1}

\begin{longtable}{|L{0.18\textwidth}|L{0.27\textwidth}|L{0.49\textwidth}|}
\hline
\textbf{Từ viết tắt} & \textbf{Từ đầy đủ} & \textbf{Giải thích} \\
\hline
\endfirsthead

\hline
\textbf{Từ viết tắt} & \textbf{Từ đầy đủ} & \textbf{Giải thích} \\
\hline
\endhead

\hline
\endfoot

AI & Artificial Intelligence & Lĩnh vực khoa học máy tính nghiên cứu các hệ thống có khả năng thực hiện tác vụ đòi hỏi trí thông minh con người. \\
\hline

LLM & Large Language Model & Mô hình học sâu với hàng tỷ tham số, có khả năng hiểu và tạo sinh ngôn ngữ tự nhiên. \\
\hline

RAG & Retrieval-Augmented Generation & Kỹ thuật kết hợp tìm kiếm thông tin từ cơ sở dữ liệu với khả năng tạo sinh của LLM để cải thiện độ chính xác. \\
\hline

API & Application Programming Interface & Tập hợp quy tắc và giao thức cho phép các phần mềm giao tiếp với nhau. \\
\hline

REST & Representational State Transfer & Kiểu kiến trúc thiết kế API trên nền giao thức HTTP, dùng cho giao tiếp giữa giao diện web và máy chủ. \\
\hline

SSR & Server-Side Rendering & Kỹ thuật kết xuất trang web phía máy chủ trước khi gửi về trình duyệt, giúp tối ưu tốc độ hiển thị lần đầu và SEO. \\
\hline

SEO & Search Engine Optimization & Tối ưu hóa để website thân thiện và dễ được các công cụ tìm kiếm xếp hạng cao. \\
\hline

SSE & Server-Sent Events & Cơ chế cho phép máy chủ đẩy dữ liệu liên tục về trình duyệt theo một chiều; được dùng cho luồng trả lời thời gian thực của trợ lý AI. \\
\hline

JWT & JSON Web Token & Chuẩn mã thông báo được ký số dùng để xác thực và truyền thông tin phiên đăng nhập một cách an toàn. \\
\hline

OAuth & Open Authorization & Giao thức ủy quyền cho phép đăng nhập thông qua tài khoản của bên thứ ba (ví dụ Google) mà không chia sẻ mật khẩu. \\
\hline

UI & User Interface & Giao diện người dùng -- phần người dùng nhìn thấy và tương tác trực tiếp với hệ thống. \\
\hline

WebSocket & WebSocket & Giao thức truyền dữ liệu hai chiều theo thời gian thực giữa trình duyệt và máy chủ trên một kết nối bền; dùng cho thông báo tức thời. \\
\hline

STOMP & Simple Text Oriented Messaging Protocol & Giao thức nhắn tin dạng văn bản chạy bên trên WebSocket, dùng để đẩy thông báo thời gian thực tới ứng dụng web. \\
\hline

SockJS & SockJS & Thư viện cung cấp kết nối kiểu WebSocket kèm cơ chế dự phòng (hỏi vòng) khi trình duyệt hoặc mạng không hỗ trợ WebSocket. \\
\hline

ICU & International Components for Unicode & Bộ định dạng chuỗi đa ngôn ngữ (số nhiều, tham số, ngày giờ) giúp hiển thị đúng ngữ pháp khi dịch giữa các ngôn ngữ. \\
\hline

NLP & Natural Language Processing & Xử lý ngôn ngữ tự nhiên -- một lĩnh vực thuộc trí tuệ nhân tạo nghiên cứu sự tương tác giữa máy tính và ngôn ngữ tự nhiên của con người. \\
\hline

RBAC & Role-Based Access Control & Kiểm soát truy cập dựa trên vai trò -- phương thức phân quyền và quản lý quyền hạn người dùng dựa trên vai trò cụ thể trong hệ thống. \\
\hline

PWA & Progressive Web Application & Ứng dụng web tiến trình -- giải pháp xây dựng ứng dụng web mang lại trải nghiệm và hiệu năng mượt mà tương tự ứng dụng di động gốc. \\
\hline

SOA & Service-Oriented Architecture & Kiến trúc hướng dịch vụ -- kiểu thiết kế phần mềm trong đó các thành phần ứng dụng cung cấp dịch vụ cho các thành phần khác thông qua giao thức truyền thông. \\
\hline

HTTP/HTTPS & HyperText Transfer Protocol (Secure) & Giao thức truyền tải dữ liệu trên nền web; phiên bản HTTPS bổ sung lớp mã hóa để bảo mật đường truyền. \\
\hline

JSON & JavaScript Object Notation & Định dạng văn bản nhẹ để biểu diễn và trao đổi dữ liệu có cấu trúc giữa các thành phần của hệ thống. \\
\hline

CI/CD & Continuous Integration / Continuous Deployment & Quy trình tự động tích hợp, kiểm thử và triển khai mã nguồn liên tục mỗi khi có thay đổi. \\
\hline

DTO & Data Transfer Object & Đối tượng truyền tải dữ liệu giữa các tầng, tách biệt cấu trúc dữ liệu trao đổi khỏi mô hình nghiệp vụ bên trong. \\
\hline

Modular Monolith & Modular Monolith & Kiến trúc nguyên khối mô-đun: ứng dụng chạy như một khối duy nhất nhưng được chia tách rõ ràng thành các module nghiệp vụ độc lập, sẵn sàng tách thành dịch vụ nhỏ khi cần. \\
\hline

Clean Architecture & Clean Architecture & Kiến trúc sạch: tổ chức mã nguồn theo nhiều tầng nhằm tách biệt logic nghiệp vụ khỏi các framework và chi tiết hạ tầng bên ngoài. \\
\hline

ReAct & Reasoning and Acting & Mô hình vận hành tác tử AI trong đó mô hình ngôn ngữ luân phiên suy luận và gọi công cụ cho tới khi đủ dữ kiện để trả lời. \\
\hline

ICL & In-Context Learning & Phương pháp học ngữ cảnh: cung cấp các ví dụ hoặc dữ liệu tham chiếu ngay trong câu lệnh (prompt) để mô hình ngôn ngữ suy luận mà không cần huấn luyện lại. \\
\hline

MCP & Model Context Protocol & Giao thức chuẩn hóa việc cung cấp ngữ cảnh và công cụ cho mô hình ngôn ngữ, cho phép các trợ lý AI bên ngoài truy vấn dữ liệu hệ thống. \\
\hline

CBF & Content-Based Filtering & Lọc theo nội dung: thuật toán gợi ý dựa trên độ tương đồng đặc trưng giữa các tin đăng. \\
\hline

CF & Collaborative Filtering & Lọc cộng tác: thuật toán gợi ý dựa trên hành vi tương tác của những người dùng tương tự nhau. \\
\hline

OTP & One-Time Password & Mã xác thực dùng một lần, có hiệu lực trong thời gian ngắn, dùng để xác minh danh tính người dùng. \\
\hline

Scrum & Scrum & Khung quản lý dự án theo phương pháp Agile, tổ chức công việc thành các chu kỳ (sprint) ngắn kèm họp định kỳ để tổng kết và lập kế hoạch. \\
\hline

% TODO: Bổ sung thêm các từ viết tắt và thuật ngữ khác khi hoàn thiện nội dung báo cáo.

\end{longtable}
